% GENERATED FILE -- do not edit by hand.
% Source: research\output\runs\defence-models-v1\result.json
% SHA-256: 63d9708d776e675b5505e65e3028e568433577538b38bcdbddb3e66591bd7c28
% Run: python scripts/generate_flmsec_tables.py
\begin{table}[t]
\centering
\small
\begin{tabular}{lccc}
\toprule
\textbf{Defence model} & \textbf{Native property Q} & \textbf{PE property} &
\textbf{Counterexample} \\
\midrule
Dual-LLM & Safe & Unsafe & 4-step PE witness \\
CaMeL & Safe & Unsafe & 4-step PE witness \\
Progent & Safe & Unsafe & 3-step PE witness \\
PACT & Safe & Unsafe & 3-step PE witness \\
Requester-only & --- & Unsafe & 4-step PE witness \\
\textbf{ITES} & Safe (PE is native) & Safe & --- \\
\bottomrule
\end{tabular}
\caption{Comparative defence verification. Each defence
abstraction satisfies its own intended property Q; a PE counterexample
demonstrates non-implication between the defence's native objective and the
Conflux PE property, not a defect in the compared system. ITES preserves PE.
All results are preliminary finite-model comparisons, not implementation-level
evaluations.}
\label{tab:comparative-defence}
\end{table}
